\documentclass[../../main.tex]{subfiles}
\graphicspath{{../../image/}{../../figure/}}
\begin{document}

\section{分层法}\label{section:ch06-分层法}

\eqref{align:ch06-6-1-1}、\eqref{align:ch06-6-1-2}的可容许弱解在此将被确定为函数族 $\{u_{h}\}$ 当 $h\downarrow 0$ 时的极限；该函数族由以分层模式拼接\eqref{align:ch06-6-1-1}的经典解而构造。除提供另一种构造解的方法并从而给出存在性定理\cref{theorem:ch06-6-2-2}的另一证明外，本途径还为可容许性条件\cref{definition:ch06-6-2-1}提供了不同的依据。

初值 $u_{0}\in L^{\infty}(\mathbb{R}^{m})$，取值于紧区间 $[a,b]$。近似解的构造涉及用核 $\lambda_{h}$ 对 $\mathbb{R}^{m}$ 上函数作磨光，$\lambda_{h}$ 构造如下。取 $\mathbb{R}$ 上非负光滑函数 $\rho$，支集含于 $[-1,1]$，偶函数，$\rho(-\xi)=\rho(\xi)$，且总质量为一（见\eqref{align:ch06-6-2-15}）。对 $h>0$，令
\begin{align}
\lambda_{h}(x)=(ph)^{-m}\prod_{\beta=1}^{m}\rho\Bigl{(}\frac{x_{\beta}}{ph}\Bigr{)},
\label{align:ch06-6-6-5-1}
\end{align}
其中
\begin{align}
p=\sqrt{m}\,q\gamma\|u_{0}\|_{L^{\infty}(\mathbb{R}^{m})},
\label{align:ch06-6-6-5-2}
\end{align}
$q$ 表示函数 $\rho$ 的总变差，$\gamma$ 是 $|G^{\prime\prime}(u)|$ 在区间 $[a,b]$ 上的最大值。用 $\lambda_{h}$ 磨光 $f\in L^{\infty}(\mathbb{R}^{m})$：
\begin{align}
(\lambda_{h}*f)(x)=\int_{\mathbb{R}^{m}}\lambda_{h}(x-y)f(y)\,\mathrm{d}y,\qquad x\in\mathbb{R}^{m}.
\label{align:ch06-6-6-5-3}
\end{align}
由\eqref{align:ch06-6-6-5-3}与\eqref{align:ch06-6-6-5-1}易见
\begin{align}
\inf(\lambda_{h}*f)\geqslant\operatorname*{ess\,inf}f,\qquad
\sup(\lambda_{h}*f)\leqslant\operatorname*{ess\,sup}f,
\label{align:ch06-6-6-5-4}
\end{align}
\begin{align}
\|\lambda_{h}*f\|_{L^{1}(\mathscr{B}_{r})}\leqslant\|f\|_{L^{1}(\mathscr{B}_{r+\sqrt{m}ph})},\qquad\text{对任意 }r>0,
\label{align:ch06-6-6-5-5}
\end{align}
\begin{align}
\|\partial_{\alpha}(\lambda_{h}*f)\|_{L^{\infty}(\mathbb{R}^{m})}\leqslant\frac{q}{ph}\|f\|_{L^{\infty}(\mathbb{R}^{m})},\qquad \alpha=1,\cdots,m.
\label{align:ch06-6-6-5-6}
\end{align}
一个稍精细的估计关键地依赖于 $\lambda_{h}$ 为偶函数，其证明见\cref{section:ch06-注记}所引文献：

\begin{align}
\Bigl{|}\int_{\mathbb{R}^{m}}\chi(x)[(\lambda_{h}*f)(x)-f(x)]\,\mathrm{d}x\Bigr{|}\leqslant ch^{2}\|\chi\|_{C^{2}(\mathbb{R}^{m})}\|f\|_{L^{\infty}(\mathbb{R}^{m})},
\label{align:ch06-6-6-5-7}
\end{align}
对一切 $\chi\in C_{0}^{\infty}(\mathbb{R}^{m})$ 成立。

近似解的构造如下。固定参数 $h>0$，把 $\mathbb{R}^{m}\times[0,\infty)$ 分层：
\begin{align}
\mathbb{R}^{m}\times[0,\infty)=\bigcup_{\ell=0}^{\infty}\mathbb{R}^{m}\times[\ell h,\ell h+h).
\label{align:ch06-6-6-5-8}
\end{align}
初值 $u_{h}(\cdot,0)$ 由
\begin{align}
u_{h}(\cdot,0)=\lambda_{h}*u_{0}(\cdot)
\label{align:ch06-6-6-5-9}
\end{align}
确定。由\eqref{align:ch06-6-6-5-6}与\eqref{align:ch06-6-6-5-2}，$u_{h}(\cdot,0)$ Lipschitz 连续，Lipschitz 常数为 $\omega=1/(p\gamma)$。于是由\cref{theorem:ch06-6-1-1}，以 $u_{h}(\cdot,0)$ 为初值的\eqref{align:ch06-6-1-1} 在层 $\mathbb{R}^{m}\times[0,h)$ 上存在经典解 $u_{h}$。

接着通过磨光 $u_{h}(\cdot,t)$ 当 $t\uparrow h$ 时的极限 $u_{h}(\cdot,h-)$ 来确定 $u_{h}(\cdot,h)$：
\begin{align}
u_{h}(\cdot,h)=\lambda_{h}*u_{h}(\cdot,h-).
\label{align:ch06-6-6-5-10}
\end{align}
在层 $\mathbb{R}^{m}\times[h,2h)$ 上，以 $t=h$ 时刻的数据 $u_{h}(\cdot,h)$ 求解\eqref{align:ch06-6-1-1}，从而把 $u_{h}$ 延拓。

继续此过程，在一般层 $[\ell h,\ell h+h)$ 上以 $t=\ell h$ 时刻的数据
\begin{align}
u_{h}(\cdot,\ell h)=\lambda_{h}*u_{h}(\cdot,\ell h-)
\label{align:ch06-6-6-5-11}
\end{align}
求解\eqref{align:ch06-6-1-1}。最终得到 $\mathbb{R}^{m}\times[0,\infty)$ 上的可测函数 $u_{h}$，取值于区间 $[a,b]$。在每层 $\mathbb{R}^{m}\times[\ell h,\ell h+h)$ 内部，$u_{h}$ 是\eqref{align:ch06-6-1-1}的经典解。然而，当跨过相邻层之间的边界 $t=\ell h$ 时，$u_{h}$ 从 $u_{h}(\cdot,\ell h-)$ 到 $u_{h}(\cdot,\ell h)$ 经历跳跃间断。

\begin{theorem}\label{theorem:ch06-6-5-1}
当 $h\downarrow 0$ 时，上述构造的族 $\{u_{h}\}$ 在 $\mathbb{R}^{m}\times[0,\infty)$ 上一致有界且几乎处处收敛于\eqref{align:ch06-6-1-1}、\eqref{align:ch06-6-1-2}的可容许解 $u$。
\end{theorem}

证明是下述两个命题与可容许解唯一性（\cref{corollary:ch06-6-2-4}）的直接推论。经典解的极限产生可容许弱解这一事实，为\cref{definition:ch06-6-2-1}提供了另一依据。

\begin{lemma}\label{lemma:ch06-6-5-2}
（相容性）假设存在序列 $\{h_{k}\}$，$h_{k}\to 0$（$k\to\infty$），使
\begin{align}
u_{h_{k}}(x,t)\to u(x,t),\qquad\text{a.e.\ 在 }\mathbb{R}^{m}\times[0,\infty)\text{ 上}.
\label{align:ch06-6-6-5-12}
\end{align}
则 $u$ 是\eqref{align:ch06-6-1-1}、\eqref{align:ch06-6-1-2}的可容许弱解。
\end{lemma}

\begin{proof}
任取凸熵函数 $\eta$，相应熵流 $Q$ 由\eqref{align:ch06-6-2-1}确定。在每层内部，$u_{h}$ 是\eqref{align:ch06-6-1-1}的经典解，故满足恒等式
\begin{align}
\partial_{t}\eta(u_{h}(x,t))+\operatorname{div}Q(u_{h}(x,t))=0.
\label{align:ch06-6-6-5-13}
\end{align}
固定 $\mathbb{R}^{m}\times[0,\infty)$ 上任一具紧支集的非负光滑试验函数 $\psi$。以 $\psi$ 乘\eqref{align:ch06-6-6-5-13}，在每层上积分，分部积分，再把所得方程对所有层求和，得
\begin{align}
&\int_{0}^{\infty}\int_{\mathbb{R}^{m}}\Bigl{[}\partial_{t}\psi\,\eta(u_{h})+\sum_{\alpha=1}^{m}\partial_{\alpha}\psi\,Q_{\alpha}(u_{h})\Bigr{]}\,\mathrm{d}x\,\mathrm{d}t+\int_{\mathbb{R}^{m}}\psi(x,0)\eta(u_{h}(x,0))\,\mathrm{d}x\notag\\
&\qquad=-\sum_{\ell=1}^{\infty}\int_{\mathbb{R}^{m}}\psi(x,\ell h)[\eta(u_{h}(x,\ell h))-\eta(u_{h}(x,\ell h-))]\,\mathrm{d}x.
\label{align:ch06-6-6-5-14}
\end{align}
结合\eqref{align:ch06-6-6-5-11}与 Jensen 不等式，并利用\eqref{align:ch06-6-6-5-7}，得
\begin{align}
&\int_{\mathbb{R}^{m}}\psi(x,\ell h)[\eta(u_{h}(x,\ell h))-\eta(u_{h}(x,\ell h-))]\,\mathrm{d}x\notag\\
&\qquad\leqslant\int_{\mathbb{R}^{m}}\psi(x,\ell h)[(\lambda_{h}*\eta(u_{h}))(x,\ell h-)-\eta(u_{h}(x,\ell h-))]\,\mathrm{d}x\leqslant Ch^{2}.
\label{align:ch06-6-6-5-15}
\end{align}
\eqref{align:ch06-6-6-5-14}右端求和含 $O(1/h)$ 个非零项。因此沿序列 $\{h_{k}\}$ 在\eqref{align:ch06-6-6-5-14}中令 $k\to\infty$，并利用\eqref{align:ch06-6-6-5-12}、\eqref{align:ch06-6-6-5-9}与\eqref{align:ch06-6-6-5-15}，可知 $u$ 满足\eqref{align:ch06-6-2-3}。证明完成。
\end{proof}

\begin{lemma}\label{lemma:ch06-6-5-3}
（紧性）存在序列 $\{h_{k}\}$，$h_{k}\to 0$（$k\to\infty$），以及 $\mathbb{R}^{m}\times[0,\infty)$ 上的 L$^{\infty}$ 函数 $u$，使\eqref{align:ch06-6-6-5-12}成立。
\end{lemma}

\begin{proof}
第一步是证明较弱的论断：存在序列 $\{h_{k}\}$（$h_{k}\to 0$，$k\to\infty$）及函数 $u$，使对几乎所有 $t\in[0,\infty)$，
\begin{align}
u_{h_{k}}(\cdot,t)\to u(\cdot,t),\qquad\text{当 }k\to\infty,\quad\text{在 }L^{\infty}(\mathbb{R}^{m})\text{ 弱$^{*}$拓扑下}.
\label{align:ch06-6-6-5-16}
\end{align}
为此，固定 $\mathbb{R}^{m}$ 上任一具紧支集的光滑试验函数 $\chi$，考虑函数
\begin{align}
v_{h}(t)=\int_{\mathbb{R}^{m}}\chi(x)u_{h}(x,t)\,\mathrm{d}x,\qquad t\in[0,\infty).
\label{align:ch06-6-6-5-17}
\end{align}
注意 $v_{h}$ 在 $[\ell h,\ell h+h)$ 上光滑，且
\begin{align}
\int_{\ell h}^{\ell h+h}\Bigl{|}\frac{\mathrm{d}}{\mathrm{d}t}v_{h}(t)\Bigr{|}\,\mathrm{d}t
&=\int_{\ell h}^{\ell h+h}\Bigl{|}-\int_{\mathbb{R}^{m}}\chi(x)\sum_{\alpha=1}^{m}\partial_{\alpha}G_{\alpha}(u(x,t))\,\mathrm{d}x\Bigr{|}\,\mathrm{d}t\notag\\
&=\int_{\ell h}^{\ell h+h}\Bigl{|}\int_{\mathbb{R}^{m}}\sum_{\alpha=1}^{m}\partial_{\alpha}\chi(x)\,G_{\alpha}(u(x,t))\,\mathrm{d}x\Bigr{|}\,\mathrm{d}t\leqslant Ch.
\label{align:ch06-6-6-5-18}
\end{align}
另一方面，$v_{h}$ 在点 $t=\ell h$ 处经历跳跃间断，可借\eqref{align:ch06-6-6-5-11}与\eqref{align:ch06-6-6-5-7}估计：
\begin{align}
|v_{h}(\ell h)-v_{h}(\ell h-)|=\Bigl{|}\int_{\mathbb{R}^{m}}\chi(x)[u_{h}(x,\ell h)-u_{h}(x,\ell h-)]\,\mathrm{d}x\Bigr{|}\leqslant Ch^{2}.
\label{align:ch06-6-6-5-19}
\end{align}
由\eqref{align:ch06-6-6-5-18}与\eqref{align:ch06-6-6-5-19}，$v_{h}$ 在 $[0,\infty)$ 的任一紧子区间上的全变差关于 $h$ 一致有界。于是由 Helly 定理（参见\cref{section:ch01-BV函数理论概述}），存在序列 $\{h_{k}\}$（$h_{k}\to 0$，$k\to\infty$），使 $v_{h_{k}}(t)$ 对几乎所有 $t\in[0,\infty)$ 收敛。


通过 Cantor 对角线程序，可构造 $\{h_{k}\}$ 的子列（仍记为 $\{h_{k}\}$），使序列
\begin{align}
\Bigl{\{}\int_{\mathbb{R}^{m}}\chi(x)u_{h_{k}}(x,t)\,\mathrm{d}x\Bigr{\}}
\label{align:ch06-6-6-5-20}
\end{align}
对几乎所有 $t\in[0,\infty)$ 收敛。
\end{proof}

\end{document}
